\chapter{2010年湖南卷（理）}

\section{单选题}

\begin{problem}
已知集合 \(M = \{1, 2, 3\}\)，\(N = \{2, 3, 4\}\)，则（\quad）

\choices
    {\(M \subseteq N\)}
    {\(N \subseteq M\)}
    {\(M \cap N = \{2, 3\}\)}
    {\(M \cup N = \{1, 4\}\)}
\end{problem}

\begin{answer}
C
\end{answer}

\begin{problem}
下列命题中是假命题的是（\quad）

\choices
    {\(\forall x\in \R,2^{x - 1} > 0\)}
    {\(\forall x\in\N^{*}\)，\((x - 1)^{2} > 0\)}
    {\(\exists x\in \R\)，\(\lg x < 1\)}
    {\(\exists x\in \R\)，\(\tan x = 2\)}
\end{problem}

\begin{answer}
B
\end{answer}

\begin{problem}
极坐标方程 \(\rho = \cos \theta\) 和参数方程 \(\begin{cases}
x = -1 - t,\\
y = 2 + 3t,
\end{cases}\)（\(t\) 为参数）所表示的图形分别是（\quad）

\choices
    {圆，直线}
    {直线，圆}
    {圆，圆}
    {直线，直线}
\end{problem}

\begin{answer}
A
\end{answer}

\begin{problem}
在 \(\mathrm{Rt}\triangle ABC\) 中，\(\angle C = 90^{\circ}\)，\(AC = 4\)，则 \( \overrightarrow{AB} \cdot \overrightarrow{AC} \) 等于（\quad）

\choices
    {\(-16\)}
    {\(-8\)}
    {8}
    {16}
\end{problem}

\begin{answer}
D
\end{answer}

\begin{problem}
\(\int_2^4 \frac{1}{x} \mathrm{d}x\) 等于（\quad）

\choices
    {\(-2\ln 2\)}
    {\(2\ln 2\)}
    {\(-\ln 2\)}
    {\(\ln 2\)}
\end{problem}

\begin{answer}
D
\end{answer}

\begin{problem}
在 \(\triangle ABC\) 中，角 \(A\)，\(B\)，\(C\) 所对的边长分别为 \(a\)，\(b\)，\(c\)，若 \(\angle C = 120^{\circ}\)，\(c = \sqrt{2} a\)，则（\quad）

\choices
    {\(a > b\)}
    {\(a < b\)}
    {\(a = b\)}
    {\(a\) 与 \(b\) 的大小关系不能确定}
\end{problem}

\begin{answer}
A
\end{answer}

\begin{problem}
在某种信息传输过程中，用 4 个数字的一个排列（数字允许重复）表示一个信息，不同排列表示不同信息. 若所用数字只有 0 和 1，则与信息 0110 至多有两个对应位置上的数字相同的信息个数为（\quad）

\choices
    {10}
    {11}
    {12}
    {15}
\end{problem}

\begin{answer}
B
\end{answer}

\begin{problem}
用 \(\min \{a, b\}\) 表示 \(a\)，\(b\) 两数中的最小值. 若函数 \(f(x) = \min \{|x|, |x + t|\}\) 的图象关于直线 \(x = -\frac{1}{2}\) 对称，则 \(t\) 的值为（\quad）

\choices
    {\(-2\)}
    {2}
    {\(-1\)}
    {1}
\end{problem}

\begin{answer}
D
\end{answer}

\section{填空题}

\begin{problem}
已知一种材料的最佳加入量在 \(110 \mathrm{~g}\) 到 \(210 \mathrm{~g}\) 之间. 若用 \(0.618\) 法安排实验，则第一次试点的加入量可以是\fillinblank{}\(\mathrm{g}\).
\end{problem}

\begin{answer}
\(171.8\) 或 \(148.2\)
\end{answer}

\begin{problem}
如图所示，过 \(\odot O\) 外一点 \(P\) 作一条直线与 \(\odot O\) 交于 \(A\)，\(B\) 两点. 已知 \(PA = 2\)，点 \(P\) 到 \(\odot O\) 的切线长 \(PT = 4\)，则弦 \(AB\) 的长为\fillinblank{}.

\bitmapfigure[max width=0.32\linewidth]{img/2010/hunan_science/q10_fig1.png}
\end{problem}

\begin{answer}
6
\end{answer}

\begin{problem}
在区间 \([-1, 2]\) 上随机取一个数 \(x\)，则 \(|x| \leqslant 1\) 的概率为\fillinblank{}.
\end{problem}

\begin{answer}
\(\frac{2}{3}\)
\end{answer}

\begin{problem}
如图是求 \(1^{2} + 2^{2} + 3^{2} + \cdots + 100^{2}\) 的值的程序框图，则正整数 \(n =\)\fillinblank{}.

\texfigure[max width=0.42\linewidth]{img_repaint/2010/hunan_science/q12_flowchart.tex}
\end{problem}

\begin{answer}
100
\end{answer}

\begin{problem}
图中的三个直角三角形是一个体积为 \(20 \mathrm{~cm}^{3}\) 的几何体的三视图，则 \(h = \)\fillinblank{}\(\mathrm{cm}\).

\bitmapfigure[max width=0.52\linewidth]{img/2010/hunan_science/q13_fig1.png}
\end{problem}

\begin{answer}
4
\end{answer}

\begin{problem}
过抛物线 \(x^{2} = 2py\) （\(p > 0\)）的焦点作斜率为 1 的直线与该抛物线交于 \(A\)，\(B\) 两点，\(A\)，\(B\) 在 \(x\) 轴上的正射影分别为 \(D\)，\(C\). 若梯形 \(ABCD\) 的面积为 \(12\sqrt{2}\)，则 \(p = \)\fillinblank{}.
\end{problem}

\begin{answer}
2
\end{answer}

\begin{problem}
若数列 \(\{a_n\}\) 满足：对任意的 \(n \in \N^{*}\)，只有有限个正整数 \(m\) 使得 \(a_m < n\) 成立. 记这样的 \(m\) 的个数为 \((a_n)^*\)，则得到一个新数列 \(\{(a_n)^*\}\). 例如，若数列 \(\{a_n\}\) 是 \(1\)，\(2\)，\(3\)，\(\cdots\)，\(n\)，\(\cdots\)，则数列 \(\{(a_n)^*\}\) 是 \(0\)，\(1\)，\(2\)，\(\cdots\)，\(n - 1\)，\(\cdots\). 已知对任意的 \(n \in \N^{*}\)，\(a_n = n^2\)，则 \((a_5)^* =\fillinblank{}\)，\(((a_n)^*)^* =\fillinblank{}\).
\end{problem}

\begin{answer}
2，\(n^2\)
\end{answer}

\section{解答题}

\begin{problem}
已知函数 \( f(x) = \sqrt{3}\sin 2x - 2\sin^2 x \).

\begin{enumerate}
\item 求函数 \(f(x)\) 的最大值；
\item 求函数 \(f(x)\) 的零点的集合.
\end{enumerate}
\begin{solution}
\begin{enumerate}
\item 由 \(2\sin^2x=1-\cos2x\)，得 \(f(x)=\sqrt{3}\sin2x+\cos2x-1=2\sin\left(2x+\frac{\pi}{6}\right)-1\). 因为 \(-1\leqslant\sin\left(2x+\frac{\pi}{6}\right)\leqslant1\)，所以 \(f(x)\leqslant1\)，且当 \(2x+\frac{\pi}{6}=\frac{\pi}{2}+2k\pi\) 时取等号，其中 \(k\in\Z\). 因此 \(f(x)\) 的最大值为 \(1\).
\item \(f(x)=0\) 等价于 \(\sin\left(2x+\frac{\pi}{6}\right)=\frac{1}{2}\). 所以 \(2x+\frac{\pi}{6}=\frac{\pi}{6}+2k\pi\) 或 \(2x+\frac{\pi}{6}=\frac{5\pi}{6}+2k\pi\)，其中 \(k\in\Z\). 解得 \(x=k\pi\) 或 \(x=k\pi+\frac{\pi}{3}\). 故零点集合为 \(\left\{k\pi\mid k\in\Z\right\}\cup\left\{k\pi+\frac{\pi}{3}\mid k\in\Z\right\}\).
\end{enumerate}
\end{solution}
\end{problem}

\begin{problem}
如图，是某城市通过抽样得到的居民某年的月均用水量（单位：吨）的频率分布直方图.

\begin{enumerate}
\item 求直方图中 \(x\) 的值；
\item 若将频率视为概率，从这个城市随机抽取 \(3\) 位居民（看作有放回的抽样），求月均用水量在 \(3\) 至 \(4\text{ 吨}\) 的居民数 \(X\) 的分布列和数学期望.
\end{enumerate}

\texfigure[max width=0.42\linewidth]{img_repaint/2010/hunan_science/q17_fig1.tex}
\begin{solution}
\begin{enumerate}
\item 由频率分布直方图中各组组距均为 \(1\)，各组频率之和为 \(1\)，得
\[
x+0.37+0.39+0.10+0.02=1,
\]
所以 \(x=0.12\).
\item 居民月均用水量在 \(3\) 至 \(4\text{ 吨}\) 的概率为 \(0.10\). 由于是有放回抽样，\(3\) 位居民是否属于该类相互独立，所以 \(X\sim B(3,0.10)\). 对 \(k=0,1,2,3\)，有
\[
P(X=k)=\mathrm C_{3}^{k}(0.10)^k(0.90)^{3-k}.
\]
因此 \(X\) 的分布列为
\begin{center}
\begin{tabular}{c|cccc}
\(X\) & \(0\) & \(1\) & \(2\) & \(3\) \\
\hline
\(P\) & \(0.729\) & \(0.243\) & \(0.027\) & \(0.001\)
\end{tabular}
\end{center}
数学期望为 \(E(X)=3\times0.10=0.3\).
\end{enumerate}
\end{solution}
\end{problem}

\begin{problem}
如图所示，在正方体 \(ABCD - A_{1}B_{1}C_{1}D_{1}\) 中，\(E\) 是棱 \(DD_{1}\) 的中点.

\begin{enumerate}
\item 求直线 \(BE\) 和平面 \(ABB_{1}A_{1}\) 所成的角的正弦值；
\item 在棱 \(C_{1}D_{1}\) 上是否存在一点 \(F\)，使 \(B_{1}F \parallel\) 平面 \(A_{1}BE\)？证明你的结论.
\end{enumerate}

\bitmapfigure[max width=0.38\linewidth]{img/2010/hunan_science/q18_fig1.png}
\begin{solution}
设正方体的棱长为 \(a\)，以 \(B\) 为原点，以 \(BC\)，\(BA\)，\(BB_{1}\) 所在直线分别为三个坐标轴，则
\begin{enumerate}
\item 于是 \(B=\left(0,0,0\right)\)，\(C=\left(a,0,0\right)\)，\(A=\left(0,a,0\right)\)，\(B_{1}=\left(0,0,a\right)\)，\(A_{1}=\left(0,a,a\right)\)，\(D=\left(a,a,0\right)\)，\(E=\left(a,a,\frac{a}{2}\right)\). 直线 \(BE\) 的方向向量可取 \(\bs{v}=\left(a,a,\frac{a}{2}\right)\). 平面 \(ABB_{1}A_{1}\) 的法向为 \(x\) 轴方向，因此直线 \(BE\) 与该平面所成的角为 \(\theta\) 时，
\[
\sin\theta=\frac{|a|}{\sqrt{a^2+a^2+\left(\frac{a}{2}\right)^2}}=\frac{2}{3}.
\]
\item 设 \(F\) 的坐标为 \(\left(a,t,a\right)\)，其中 \(0\leqslant t\leqslant a\). 有 \(\overrightarrow{BA_{1}}=\left(0,a,a\right)\)，\(\overrightarrow{BE}=\left(a,a,\frac{a}{2}\right)\)，\(\overrightarrow{B_{1}F}=\left(a,t,0\right)\).
取平面 \(A_{1}BE\) 的一个法向量
\[
\bs{n}=\overrightarrow{BA_{1}}\times\overrightarrow{BE}=a^2\left(-\frac{1}{2},1,-1\right).
\]
要使 \(B_{1}F\parallel\) 平面 \(A_{1}BE\)，必须且只需 \(\overrightarrow{B_{1}F}\cdot\bs{n}=0\)，即
\[
-\frac{a^3}{2}+a^2t=0,
\]
从而 \(t=\frac{a}{2}\). 因此存在这样的点 \(F\)，且 \(F\) 是棱 \(C_{1}D_{1}\) 的中点.
\end{enumerate}
\end{solution}
\end{problem}

\begin{problem}
为了考察冰川的融化状况，一支科考队在某冰川上相距 \(8 \mathrm{~km}\) 的 \(A\)，\(B\) 两点各建一个考察基地. 视冰川面为平面形，以过 \(A\)，\(B\) 两点的直线为 \(x\) 轴，线段 \(AB\) 的垂直平分线为 \(y\) 轴建立平面直角坐标系（如图）. 在直线 \(x = 2\) 的右侧，考察范围为到点 \(B\) 的距离不超过 \(\frac{6\sqrt{5}}{5} \mathrm{~km}\) 的区域. 在直线 \(x = 2\) 的左侧，考察范围为到 \(A\)，\(B\) 两点的距离之和不超过 \(4\sqrt{5} \mathrm{~km}\) 的区域.

\begin{enumerate}
\item 求考察区域边界曲线的方程；
\item 如图所示，设线段 \(P_{1}P_{2}\)，\(P_{2}P_{3}\) 是冰川的部分边界线（不考虑其他边界），当冰川融化时，边界线沿与其垂直的方向朝考察区域平行移动，第一年移动 \(0.2 \mathrm{~km}\)，以后每年移动的距离为前一年的 2 倍. 求冰川边界线移动到考察区域所需的最短时间.
\end{enumerate}

\bitmapfigure[max width=0.46\linewidth]{img/2010/hunan_science/q19_fig1.png}
\begin{solution}
\begin{enumerate}
\item 在直线 \(x=2\) 的右侧，考察区域是以 \(B\left(4,0\right)\) 为圆心、半径为 \(\frac{6\sqrt{5}}{5}\) 的圆的内部，因此其边界曲线为
\[
C_{1}:\left(x-4\right)^2+y^2=\frac{36}{5},\qquad x\geqslant2.
\]
在直线 \(x=2\) 的左侧，考察区域边界是以 \(A\left(-4,0\right)\)，\(B\left(4,0\right)\) 为焦点的椭圆. 设其长半轴为 \(\alpha\)，半焦距为 \(c_{0}\)，则 \(2\alpha=4\sqrt{5}\)，\(c_{0}=4\)，从而 \(\alpha=2\sqrt{5}\)，短半轴平方为 \(\alpha^2-c_{0}^2=4\). 所以另一段边界曲线为
\[
C_{2}:\frac{x^2}{20}+\frac{y^2}{4}=1,\qquad x<2.
\]
\item 直线段 \(P_{1}P_{2}\) 的斜率为
\[
\frac{6-(-1)}{-\frac{8\sqrt{3}}{3}+5\sqrt{3}}=\sqrt{3},
\]
故其所在直线方程为 \(\sqrt{3}x-y+14=0\). 对椭圆内部的点，有
\[
\sqrt{3}x-y=2\left(\sqrt{15}\frac{x}{2\sqrt{5}}-\frac{y}{2}\right)\geqslant-2\sqrt{15+1}=-8.
\]
等号在 \(\left(-\frac{5\sqrt{3}}{2},\frac{1}{2}\right)\) 处取得，且该点满足 \(x<2\)，属于考察区域边界. 所以从直线 \(P_{1}P_{2}\) 沿法向朝考察区域移动，首次接触考察区域所需的距离为
\[
\frac{14-8}{\sqrt{3^2+(-1)^2}}=3\text{ km}.
\]
直线段 \(P_{2}P_{3}\) 所在直线为 \(y=6\). 考察区域的最高点纵坐标为圆 \(C_{1}\) 的最高点纵坐标 \(\frac{6\sqrt{5}}{5}\)，所以该线段向下移动至考察区域至少需要
\[
6-\frac{6\sqrt{5}}{5}>3\text{ km}.
\]
因此冰川边界线首次移动到考察区域时，所需移动距离为 \(3\text{ km}\).

设经过 \(n\) 年移动的总距离为 \(s_n\)，则
\[
s_n=0.2\left(1+2+2^2+\cdots+2^{n-1}\right)=0.2\left(2^n-1\right).
\]
有 \(s_3=1.4<3\)，\(s_4=3\)，故冰川边界线移动到考察区域所需的最短时间为 \(4\) 年.
\end{enumerate}
\end{solution}
\end{problem}

\begin{problem}
已知函数 \( f(x) = x^2 + bx + c \) （\( b, c \in \R \)），对任意的 \( x \in \R \)，恒有 \( f'(x) \leqslant f(x) \).

\begin{enumerate}
\item 证明：当 \(x \geqslant 0\) 时，\(f(x) \leqslant (x + c)^2\)；
\item 若对满足题设条件的任意 \(b\)，\(c\)，不等式 \(f(c) - f(b) \leqslant M (c^2 - b^2)\) 恒成立，求 \(M\) 的最小值.
\end{enumerate}
\begin{solution}
\begin{enumerate}
\item 由题设，\(f'(x)=2x+b\)，所以对任意 \(x\in\R\)，有
\[
f(x)-f'(x)=x^2+(b-2)x+c-b\geqslant0.
\]
这是关于 \(x\) 的二次函数，故其判别式不大于 \(0\)，即
\[
(b-2)^2-4(c-b)\leqslant0.
\]
整理得 \(c\geqslant\frac{b^2}{4}+1\). 当 \(x\geqslant0\) 时，由 \(c\geqslant1\) 以及
\[
2c-b\geqslant\frac{b^2}{2}-b+2=\frac{1}{2}(b-1)^2+\frac{3}{2}>0,
\]
可得
\[
(x+c)^2-f(x)=(2c-b)x+c(c-1)\geqslant0.
\]
因此 \(f(x)\leqslant(x+c)^2\).
\item 直接计算得
\[
f(c)-f(b)=c^2+bc-2b^2.
\]
取 \(M=\frac{3}{2}\)，则
\[
\frac{3}{2}(c^2-b^2)-\bigl(f(c)-f(b)\bigr)=\frac{1}{2}(c-b)^2\geqslant0,
\]
所以 \(M=\frac{3}{2}\) 时不等式对所有满足题设条件的 \(b\)，\(c\) 恒成立.

再取 \(b=2\)，\(c=2+\varepsilon\)，其中 \(\varepsilon>0\). 此时 \(c\geqslant\frac{b^2}{4}+1\)，且
\[
\frac{f(c)-f(b)}{c^2-b^2}=\frac{c+4}{c+2}=\frac{6+\varepsilon}{4+\varepsilon}\longrightarrow\frac{3}{2}\quad\left(\varepsilon\to0^+\right).
\]
因此任何小于 \(\frac{3}{2}\) 的数都不能使题设不等式恒成立，故 \(M\) 的最小值为 \(\frac{3}{2}\).
\end{enumerate}
\end{solution}
\end{problem}

\begin{problem}
数列 \(\{a_{n}\}\) （\(n\in\N^{*}\)） 中， \(a_1 = a\)，\(a_{n + 1}\) 是函数 \(f_{n}(x) = \frac{1}{3}x^{3} - \frac{1}{2}(3a_n + n^2)x^2 + 3n^2a_nx\) 的极小值点.

\begin{enumerate}
\item 当 \(a = 0\) 时，求通项 \(a_{n}\)；
\item 是否存在 \(a\) 使数列 \(\{a_{n}\}\) 是等比数列？若存在，求 \(a\) 的取值范围；若不存在，请说明理由.
\end{enumerate}

\begin{solution}
\begin{enumerate}
\item 由
\[
f_n'(x)=x^2-(3a_n+n^2)x+3n^2a_n=(x-3a_n)(x-n^2)
\]
可知，当 \(3a_n<n^2\) 时，\(x=n^2\) 是极小值点；当 \(3a_n>n^2\) 时，\(x=3a_n\) 是极小值点. 当 \(3a_n=n^2\) 时，导函数只有二重根，函数没有极小值点. 因此在数列有定义时，
\[
a_{n+1}=\max\left\{3a_n,n^2\right\}.
\]
当 \(a=0\) 时，依次有
\[
a_1=0,\qquad a_2=1,\qquad a_3=4,\qquad a_4=12.
\]
对 \(n\geqslant3\)，有 \(4\cdot3^{n-2}>n^2\). 事实上，\(n=3\) 时不等式成立；若 \(n\geqslant3\)，则左边每次乘以 \(3\)，而右边的增长倍数为 \(\frac{(n+1)^2}{n^2}\leqslant\frac{16}{9}<3\)，所以不等式对所有 \(n\geqslant3\) 成立. 于是 \(a_{n+1}=3a_n\)，从而
\[
a_n=4\cdot3^{n-3},\qquad n\geqslant3.
\]
\item 若 \(\{a_n\}\) 是等比数列，则由递推式知 \(a_2>0\)，\(a_3>0\). 若 \(a<0\)，则 \(\frac{a_2}{a_1}<0\)，而 \(\frac{a_3}{a_2}>0\)，不可能等比；若 \(a=0\)，则第一问中的数列也不等比. 因此 \(a>0\).

当 \(0<a<\frac{1}{3}\) 时，\(a_2=1\)，\(a_3=4\). 若数列等比，则公比既为 \(\frac{a_2}{a_1}=\frac{1}{a}\)，又为 \(\frac{a_3}{a_2}=4\)，所以 \(a=\frac{1}{4}\). 但此时 \(a_4=\max\{12,9\}=12\)，而等比数列应有 \(a_4=16\)，矛盾. 当 \(a=\frac{1}{3}\) 时，\(3a_1=1^2\)，\(f_1'(x)=(x-1)^2\)，函数没有极小值点，数列不能按题意定义.

当 \(a>\frac{1}{3}\) 时，\(a_2=3a\). 若数列等比，则公比为 \(3\)，从而必须有 \(a_3=9a\). 当 \(a<\frac{4}{9}\) 时，\(a_3=\max\{9a,4\}=4\neq9a\)，矛盾；当 \(a=\frac{4}{9}\) 时，\(3a_2=4=2^2\)，此时 \(f_2'(x)=(x-4)^2\)，没有极小值点，数列不能按题意定义. 所以必要条件为 \(a>\frac{4}{9}\).

反之，若 \(a>\frac{4}{9}\)，则 \(a_2=3a\)，且 \(3a>1\). 又 \(9a>4\)，并且对 \(n\geqslant2\)，若 \(3^na>n^2\)，则
\[
3^{n+1}a>3n^2>(n+1)^2.
\]
因此归纳可得 \(3^na>n^2\) 对所有 \(n\geqslant1\) 成立，递推式始终取 \(a_{n+1}=3a_n\). 所以 \(a_n=3^{n-1}a\) 是等比数列.

综上，存在这样的 \(a\)，其取值范围为
\[
a\in\left(\frac{4}{9},+\infty\right).
\]
\end{enumerate}
\end{solution}
\end{problem}
